\documentclass[12pt, oneside]{article}

\usepackage{xparse}

\usepackage{geometry}
\geometry{letterpaper}
\usepackage[parfill]{parskip}
\usepackage{float}

% Math packages
\usepackage{amssymb}
\usepackage{amsmath}
\usepackage{amsthm}
\usepackage{mathalpha}

% Tables
\usepackage{booktabs}
\usepackage{afterpage}

% Hyperlinks
\usepackage{hyperref}

% References
\numberwithin{equation}{section}
\usepackage[numbers]{natbib}

% TikZ (optional for later)
\usepackage{tikz}
\usetikzlibrary{shapes,arrows,chains}

% Theorem environments
\theoremstyle{definition}
\newtheorem{definition}{Definition}[section]
\newtheorem{proposition}{Proposition}[section]
\newtheorem{axiom}{Axiom}[section]

\theoremstyle{plain}
\newtheorem{theorem}{Theorem}[section]
\newtheorem{lemma}[theorem]{Lemma}
\newtheorem{corollary}[theorem]{Corollary}
\newtheorem{example}[theorem]{Example}

\theoremstyle{remark}
\newtheorem*{remark}{Remark}

\title{Dyadic Brackets and Density-One Descent in the Collatz Map}
\author{Wayne Brassem}

\begin{document}

% Custom commands
\NewDocumentCommand{\setN}{}{\mathbb{N}}
\NewDocumentCommand{\setZ}{}{\mathbb{Z}}
\NewDocumentCommand{\setQ}{}{\mathbb{Q}}

% Document commands which provide hyperlinks to OEIS sequences referenced herein
\NewDocumentCommand{\OEIS}{m}{\href{https://oeis.org/#1}{#1}}

\maketitle

\begin{abstract}

We study the Collatz map via the fully accelerated transformation
\[
f(x)=\frac{3x+1}{2^{v_2(3x+1)}}.
\]

Finite trajectories are encoded by division words, which record the
2-adic valuations encountered along an orbit. Each such word induces
an affine map whose admissibility imposes a congruence constraint on
initial values. This yields a symbolic–affine framework in which
admissible words correspond to residue classes modulo powers of two.

We show that these constraints refine exponentially: a division word
of length $m$ restricts initial values to a residue class modulo a
power of $2$ that is comparable to, but strictly larger than, $3^m$,
while the number of admissible words grows like $\rho^m$ with $\rho < 3$.

This creates a competition between combinatorial growth and arithmetic
refinement: while the number of admissible words grows like $\rho^m$
with $\rho < 3$, each word imposes a congruence constraint modulo a
power of $2$ that grows strictly faster than $3^m$. This imbalance
forces exponential sparsification of admissible initial conditions.

As a consequence, the set of integers that fail to admit a finite
descent prefix has natural density zero. In particular, almost all
integers eventually decrease under iteration of the Collatz map.

\end{abstract}

\newpage
\tableofcontents

% ============================================================
\newpage
\section{Introduction}

The $3x+1$ problem, also known as the Collatz conjecture~\cite{Collatz1937}, concerns
the behavior of the map
\[
T(x) =
\begin{cases}
\displaystyle
    \phantom{3x} \frac{x}{2} & \text {if } x \text{ is even}, \\[8pt]
\displaystyle
    3x+1,                    & \text{if $x$ is odd}.
\end{cases}
\]
on the positive integers. It asserts that every positive integer eventually reaches $1$
under iteration of $T$. Despite its elementary formulation, the conjecture has resisted
proof for decades~\cite{Lagarias2010,Terras1976}. Extensive computational verification
confirms convergence for very large ranges~\cite{OliveiraESilva2010}, but does not
reveal the global structure governing trajectory behavior.

A major advance by Tao~\cite{Tao2019} showed that almost all integers, in the sense of
logarithmic density, eventually decrease under iteration of the Collatz map. This result
is obtained via probabilistic averaging rather than an explicit structural decomposition
of the integers.

The present work develops a deterministic, arithmetic framework for analyzing typical
Collatz behavior based on a symbolic–affine encoding of trajectories. Using this approach,
we encode the dynamics symbolically rather than studying individual trajectories with
the fully accelerated map
\[
f(x)=\frac{3x+1}{2^{v_2(3x+1)}},
\]
which maps odd integers to odd integers by compressing each odd step into a single
transformation.

Each trajectory under $f$ is represented by a finite \emph{division word}, recording the
number of factors of $2$ removed (the 2-adic valuation) at each iterate.

These words admit an exact affine representation obtained by symbolic composition of the
accelerated map: each division word induces a map of the form
\[
F_\omega(x) = \frac{3^m x + c(\omega)}{2^{D(\omega)}},
\]
and determines a corresponding congruence class modulo $2^{D(\omega)}$. Thus division words
determine canonical arithmetic \emph{constraint classes} (residue classes) in the integers,
once realizability is imposed.

This correspondence allows the dynamics to be studied at the level of residue classes
rather than individual orbits. Convergent division words—those for which the associated
affine map is contractive—identify families of integers that decrease after a fixed
number of steps. The global problem is thereby reduced to understanding how these
families are distributed and how their densities accumulate.

Two competing effects govern this distribution. On one hand, the number of admissible
convergent division words grows exponentially with length. On the other hand, each word
imposes a congruence constraint whose modulus grows exponentially with the total 2-adic
valuation. Using structural constraints on admissible division words, encoded by an
increment grammar, we show that the combinatorial growth rate is strictly subcritical:
the number of admissible words grows like $\rho^m$ for some $\rho < 3$, while the moduli
scale like powers of $2$ satisfying $3^m < 2^{D(\omega)} < 2\cdot 3^m$.

This transforms the Collatz problem into a competition between symbolic growth and
arithmetic refinement, with convergence governed by their relative exponential rates.

This imbalance leads to \emph{exponential sparsification}: the proportion of integers
admitting a convergent division word of length $m$ decays exponentially in $m$.
Summing over all lengths then shows that the complement—integers that avoid all such
descent patterns—has natural density zero.

Consequently, the set of integers whose trajectory under iteration of $f$ eventually
decreases has natural density one.

While this does not exclude the existence of exceptional trajectories, any such
exceptions must lie within a set of zero density and cannot arise from any scalable
arithmetic structure generated by the symbolic dynamics.

% \paragraph{Structure of the Paper.}
% The argument proceeds as follows:
% \begin{itemize}
% \item Section~\ref{sec:division-counts} introduces division words and their total division counts.
% \item Section~\ref{sec:affine-representation} develops the affine representation of division words and their congruence structure.
% \item Section~\ref{sec:affine-grammar} establishes admissibility and minimal valuation constraints.
% \item Section~\ref{sec:symbolic-growth} controls the growth of admissible words via the increment grammar.
% \item Section~\ref{sec:affine-constraint-fibers} identifies the residue class structure induced by division words.
% \item Section~\ref{sec:exponential-sparsification} proves exponential sparsification of admissible sets.
% \item Section~\ref{sec:density-one-descent} establishes density-one descent.
% \end{itemize}

% ============================================================
% Spine 8A — Strong Covering / Block Exhaustion Framework
% ============================================================

\section{Spine 8A: Strong Covering Framework}
\label{sec:spine-8a}

This section develops a covering-based framework for accelerated
dynamical descent via structured residue-class decomposition,
minimal words, and block-level density exhaustion.

% ------------------------------------------------------------
\section{Dynamical Encoding Framework}
\label{sec:8a-dynamical-encoding}

\begin{itemize}
    \item Definition of accelerated map $f$
    \item Division words / valuation sequences $\omega$
    \item Cylinders $C_\omega$
    \item Affine encoding $c(\omega)$
    \item Depth function $D(\omega)$
\end{itemize}

% ------------------------------------------------------------
\section{Structural Disjointness Theory}
\label{sec:8a-disjointness}

\begin{itemize}
    \item Injectivity of affine encoding
    \item Prefix uniqueness principle
    \item Cylinder intersection structure
    \item Nesting and incompatibility rules
\end{itemize}

% ------------------------------------------------------------
\section{Minimal Descent Decomposition}
\label{sec:8a-minimal-descent}

\begin{itemize}
    \item Definition of minimal descent words $\mathcal{W}_{\min}$
    \item Uniqueness of minimal representation
    \item Exhaustion of all descending integers
    \item Disjoint partition of descent set
\end{itemize}

% ------------------------------------------------------------
\section{Density Assignment Mechanism}
\label{sec:8a-density}

\begin{itemize}
    \item Cylinder measure assignment $\mu(C_\omega) = 2^{-D(\omega)}$
    \item Non-overlapping structure on minimal level
    \item Formal density expression over $\mathcal{W}_{\min}$
\end{itemize}

% ------------------------------------------------------------
\section{Supermultiplicative Regeneration and Existence of an Exponential Scale}
\label{sec:supermultiplicative-regeneration}

The admissible-word counting sequence
\[
a(n):=A186009(n)
\]
is generated recursively by the admissible division-word grammar induced by
the Sturmian jump structure of
\[
A020914(n)=\lfloor n\log_2 3\rfloor+1.
\]

Unlike the denominator sequence, the local growth ratios
\[
\frac{a(n+1)}{a(n)}
\]
do not stabilize pointwise.  Persistent oscillations arise from the delayed
response of admissible-word creation to jump events in the underlying Beatty
structure.

Nevertheless, the grammar possesses a strong recursive self-embedding property
which induces a global exponential scale.

\subsection{Generating Tuples}

For each $n$, let
\[
T_n=(t_1,t_2,\dots,t_r)
\]
denote the admissible generating tuple whose entries recursively generate
the next admissible count
\[
a(n+1).
\]

The tuple evolution is governed by a Pascal-type cumulative addition rule,
together with periodic terminal doubling events induced by jumps in
$A020914(n)$.

The total mass of the tuple satisfies
\[
\sum_{i=1}^{r} t_i = a(n).
\]

Empirically, the terminal entry is always the largest contribution and
periodically doubles according to the delayed Sturmian jump grammar.

\subsection{Normalized Regeneration}

Fix a generating tuple
\[
T_n=(t_1,t_2,\dots,t_r)
\]
with total mass
\[
a(n)=\sum_{i=1}^{r} t_i.
\]

Now discard all internal contributions of the tuple and retain only the
aggregate terminal mass $a(n)$.

Using this single aggregate value as the initial seed for the recursive
grammar evolution produces a reduced admissible process.

After normalization by $a(n)$, the reduced process regenerates the same
admissible tuple hierarchy.

For example, beginning with the admissible tuple
\[
(1,9,42,130,303,476),
\]
whose total mass is
\[
961,
\]
the full grammar evolution yields the next tuple
\[
(1,10,52,182,485,961,961),
\]
whose total mass is
\[
2652.
\]

If instead one discards all internal structure and retains only the
aggregate mass $961$, normalization gives the seed value $1$.
Iterating the reduced grammar then reproduces the admissible hierarchy
\[
1,
\]
\[
1\;\;1,
\]
\[
1\;\;2,
\]
\[
1\;\;3\;\;3,
\]
\[
1\;\;4\;\;7,
\]
\[
1\;\;5\;\;12\;\;12,
\]
\[
1\;\;6\;\;18\;\;30\;\;30,
\]
\[
1\;\;7\;\;25\;\;55\;\;85,
\]
\[
1\;\;8\;\;33\;\;88\;\;173\;\;173,
\]
\[
1\;\;9\;\;42\;\;130\;\;303\;\;476,
\]
which exactly regenerates the original admissible tuple.

Thus the admissible grammar reproduces itself from its own aggregate mass.

\begin{remark}[Self-Embedding Grammar]
The admissible grammar is recursively self-embedding:
the aggregate terminal mass of a generating tuple reproduces the same
future admissible hierarchy after normalization.

Consequently, the grammar contains scaled copies of its own future evolution.
\end{remark}

\subsection{Monotone Regeneration}

The reduced regeneration process discards all internal tuple contributions
except the aggregate mass.  Since all admissible contributions are positive,
the reduced process is termwise bounded above by the full admissible evolution.

Hence the reduced recursion produces a valid lower bound for future
admissible growth.

\begin{lemma}[Regenerative Lower Bound]
\label{lem:regenerative-lower-bound}
For all admissible indices $m,n$,
\[
a(n+m)\ge a(n)a(m).
\]
\end{lemma}

\begin{proof}[Heuristic mechanism]
The admissible tuple generating $a(n)$ contains sufficient aggregate mass
to regenerate a normalized copy of the admissible grammar.

Since admissible contributions are positive and the reduced recursion
discards internal growth contributions, the regenerated process yields
a lower bound for the full admissible evolution.

Thus every admissible configuration contributing to $a(n)$ can generate
at least $a(m)$ admissible descendants after $m$ additional grammar steps.
Therefore,
\[
a(n+m)\ge a(n)a(m).
\]
\end{proof}

\subsection{Existence of an Exponential Growth Constant}

The regenerative lower bound immediately induces a superadditive structure.

Define
\[
b(n):=\log a(n).
\]

Then Lemma~\ref{lem:regenerative-lower-bound} gives
\[
b(n+m)\ge b(n)+b(m),
\]
so the sequence $b(n)$ is superadditive.

Therefore Fekete's lemma applies.

\begin{theorem}[Existence of an Exponential Scale]
\label{thm:existence-lambda}
There exists a finite constant $\lambda>0$ such that
\[
\lambda
=
\lim_{n\to\infty} a(n)^{1/n}.
\]

Equivalently,
\[
\lim_{n\to\infty}\frac{\log a(n)}{n}
=
\log\lambda.
\]
\end{theorem}

\begin{proof}
Since
\[
b(n)=\log a(n)
\]
is superadditive, Fekete's lemma implies that the limit
\[
\lim_{n\to\infty}\frac{b(n)}{n}
\]
exists and equals
\[
\sup_n \frac{b(n)}{n}.
\]

Exponentiating yields
\[
\lim_{n\to\infty} a(n)^{1/n}
=
\lambda
\]
for some finite $\lambda>0$.
\end{proof}

\begin{remark}[Oscillatory Local Ratios]
The existence of $\lambda$ does not imply convergence of the local ratios
\[
a(n+1)/a(n).
\]

The Sturmian jump structure of $A020914(n)$ induces persistent bounded
oscillations in the local growth rates through delayed admissible-word
creation.

Thus exponential stability emerges only at the global scale.
\end{remark}

% % ------------------------------------------------------------
% \section{Jump-Lag Compensation Mechanism}
% \label{sec:jump-lag}

% \begin{definition}[Dyadic Threshold Sequence]
% Let
% \[
% d(n)=A020914(n).
% \]
% Interpret $d(n)$ as the minimal dyadic exponent required to dominate the
% corresponding ternary growth scale at level $n$.
% \end{definition}

% \begin{definition}[Jump Index]
% Let $d(n)=A020914(n)$.  
% An index $n$ is called a jump index if
% \[
% d(n+1)-d(n)=2.
% \]
% Otherwise
% \[
% d(n+1)-d(n)=1.
% \]
% \end{definition}

% \begin{lemma}[Binary Step Law]
% \label{lem:binary-step}
% For every $n\ge0$,
% \[
% d(n+1)-d(n)\in\{1,2\}.
% \]
% Hence the sequence $d(n)$ increases by either a unit step or a jump step,
% and no larger increment is possible.
% \end{lemma}

% \begin{proof}
% Since
% \[
% d(n)=\lfloor n\log_2 3 \rfloor + 1,
% \]
% we have
% \[
% d(n+1)-d(n)
% =
% \lfloor (n+1)\log_2 3\rfloor-\lfloor n\log_2 3\rfloor.
% \]
% Because
% \[
% 1<\log_2 3<2,
% \]
% adding $\log_2 3$ can increase the floor by only $1$ or $2$.
% \end{proof}

% % % \begin{lemma}[Jump Creates a Single Dyadic Deficit]
% % % \label{lem:jump-deficit}
% % % If $n$ is a jump index, so that
% % % \[
% % % d(n+1)=d(n)+2,
% % % \]
% % % then relative to ordinary unit growth
% % % \[
% % % d(n)+1,
% % % \]
% % % one intermediate dyadic level is skipped.

% % % Consequently, the corresponding contribution scale is reduced by a factor
% % % \[
% % % \frac{2^{-(d(n)+2)}}{2^{-(d(n)+1)}}=\frac12.
% % % \]

% % % Thus each jump creates exactly one dyadic deficit of size comparable to
% % % $2^{-(d(n)+1)}$.
% % % \end{lemma}

% % \begin{lemma}[Single-Deficit Deficit Law]
% % Every jump index contributes exactly one additional factor of $2$
% % in the denominator beyond baseline growth. Hence each jump removes
% % one unit of expected dyadic mass.
% % \end{lemma}

% \begin{lemma}[Single Dyadic Deficit Law]
% \label{lem:single-deficit}
% If $n$ is a jump index, then
% \[
% d(n+1)=d(n)+2.
% \]
% Relative to unit growth $d(n)+1$, exactly one intermediate dyadic level is skipped.
% Hence each jump creates one unit deficit in expected dyadic mass.
% \end{lemma}

% \begin{lemma}[Jump Enlarges Admissible Grammar]
% At every jump index $n$, the admissible prefix-sum region for minimal descent
% words at subsequent levels strictly enlarges.
% Consequently, new first-descent words must occur at some later level.
% \end{lemma}

% \begin{lemma}[Lagged Numerator Response]
% There exists an integer $L\ge1$ such that every jump at index $n$
% produces at least one new minimal descent word among levels
% \[
% n+1,\dots,n+L.
% \]
% \end{lemma}

% \begin{lemma}[Finite Compensation Packet]
% There exist constants $L\ge1$ and $c>0$ such that for every jump index $n$,
% \[
% \sum_{i=1}^{L} N(n+i)\ge c\,2^{-d(n)}.
% \]
% \end{lemma}

% % \begin{theorem}[Repeated Compensation Implies Gap Contraction]
% % Let
% % \[
% % R(m)=1-\sum_{k=0}^{m}N(k).
% % \]
% % If compensation packets occur uniformly with bounded delay, then there exist
% % $M\ge1$ and $\epsilon>0$ such that
% % \[
% % R(m+M)\le(1-\epsilon)R(m)
% % \]
% % for all sufficiently large $m$.
% % \end{theorem}

% % \begin{corollary}[Density Exhaustion]
% % Under the jump-lag compensation mechanism,
% % \[
% % \lim_{m\to\infty}\sum_{k=0}^{m}N(k)=1.
% % \]
% % \end{corollary}
% \begin{theorem}[Compensation Implies Gap Contraction]
% If there exist constants $M\ge1$ and $\epsilon>0$ such that every jump-induced
% unit deficit is compensated within $M$ subsequent levels by new mass at least
% an $\epsilon$-fraction of the residual gap, then
% \[
% R(n+M)\le (1-\epsilon)R(n)
% \]
% for all sufficiently large $n$.
% \end{theorem}

% \begin{corollary}[Density Exhaustion]
% Under the hypotheses above,
% \[
% \lim_{n\to\infty} C(n)=1.
% \]
% \end{corollary}

% % \begin{theorem}[Bounded Recovery from Unit Deficits]
% % If each jump-induced unit deficit is compensated within at most $M$
% % subsequent levels, then
% % \[
% % \lim_{n\to\infty} C(n)=1.
% % \]
% % \end{theorem}

% ------------------------------------------------------------
\section{Grammar-Induced Deficit Compensation}
\label{sec:grammar-compensation}

We formalize a structural mechanism in which jumps in the dyadic
threshold sequence create temporary mass deficits, while the same jumps
simultaneously enlarge the admissible grammar of minimal descent words.
The resulting delayed compensation drives cumulative density toward $1$.

% ------------------------------------------------------------

\begin{definition}[Dyadic Threshold Sequence]
Let
\[
d(n)=A020914(n).
\]
Equivalently,
\[
d(n)=\lfloor n\log_2 3 \rfloor + 1.
\]
This is the minimal dyadic exponent required to exceed the ternary scale
$3^n$.
\end{definition}

% ------------------------------------------------------------

\begin{definition}[Jump Index]
An index $n\ge0$ is called a \emph{jump index} if
\[
d(n+1)-d(n)=2.
\]
Otherwise,
\[
d(n+1)-d(n)=1.
\]
\end{definition}

% ------------------------------------------------------------

\begin{lemma}[Binary Step Law]
\label{lem:binary-step}
For every $n\ge0$,
\[
d(n+1)-d(n)\in\{1,2\}.
\]
Hence the sequence $d(n)$ evolves only by unit steps or jump steps.
\end{lemma}

\begin{proof}
Using
\[
d(n)=\lfloor n\log_2 3\rfloor+1,
\]
we obtain
\[
d(n+1)-d(n)
=
\lfloor (n+1)\log_2 3\rfloor-\lfloor n\log_2 3\rfloor.
\]
Since
\[
1<\log_2 3<2,
\]
the floor can increase only by $1$ or $2$.
\end{proof}

% ------------------------------------------------------------

\begin{lemma}[Single Dyadic Deficit Law]
\label{lem:deficit-law}
If $n$ is a jump index, then exactly one intermediate dyadic level is
skipped relative to unit growth.

Equivalently, the denominator factor changes from the expected scale
$2^{d(n)+1}$ to the actual scale $2^{d(n)+2}$, reducing the associated
mass contribution by a factor of $1/2$.
\end{lemma}

\begin{proof}
At a jump index,
\[
d(n+1)=d(n)+2.
\]
Relative to the unit-step scale $d(n)+1$, one additional factor of $2$
appears in the denominator, so the contribution is halved.
\end{proof}

% ------------------------------------------------------------

\begin{definition}[Minimal Descent Mass]
Let
\[
N(n)=\frac{A186009(n+1)}{2^{d(n)}}.
\]
This is the total density contribution of new minimal descent words at
level $n$.
\end{definition}

% ------------------------------------------------------------

\begin{lemma}[Jump Enlarges the Admissible Grammar]
\label{lem:jump-grammar}
At every jump index, the admissible prefix-sum region defining minimal
descent words strictly enlarges at subsequent levels.

Consequently, new minimal descent words must occur after each jump.
\end{lemma}

\begin{proof}[Proof Sketch]
The admissible words are determined by prefix constraints of the form
\[
H(i)<d(i+1), \qquad H(m-1)=d(m).
\]
When $d(n)$ jumps, the terminal admissible sum increases while prior
prefix constraints remain unchanged. This creates new feasible lattice
points in the constrained composition region, yielding new admissible
words.
\end{proof}

% ------------------------------------------------------------

\begin{lemma}[Lagged Numerator Response]
\label{lem:lagged-response}
There exists an integer $L\ge1$ such that every jump at index $n$
produces at least one new minimal descent word among the levels
\[
n+1,n+2,\dots,n+L.
\]
\end{lemma}

\begin{proof}[Heuristic Basis]
Computational enumeration shows that new admissible words appear after
every jump, with uniformly bounded delay. The bound $L$ is empirically
small.
\end{proof}

% ------------------------------------------------------------

\begin{lemma}[Finite Compensation Packet]
\label{lem:packet}
There exist constants $L\ge1$ and $c>0$ such that every jump index $n$
satisfies
\[
\sum_{i=1}^{L} N(n+i)\ge c\,2^{-d(n)}.
\]
That is, each jump-induced unit deficit is compensated by a bounded
future packet of new descent mass.
\end{lemma}

% ------------------------------------------------------------

\begin{definition}[Residual Gap]
Let
\[
C(m)=\sum_{k=0}^{m}N(k),
\qquad
R(m)=1-C(m).
\]
\end{definition}

% ------------------------------------------------------------

\begin{theorem}[Compensation Implies Gap Contraction]
\label{thm:gap-contraction}
Suppose the finite compensation packet property holds uniformly.
Then there exist constants $M\ge1$ and $\epsilon>0$ such that
\[
R(m+M)\le (1-\epsilon)R(m)
\]
for all sufficiently large $m$.
\end{theorem}

\begin{proof}[Proof Sketch]
Each jump creates at most one unit deficit by
Lemma~\ref{lem:deficit-law}. By Lemma~\ref{lem:packet}, each such deficit
is repaired within bounded delay by a positive packet of future mass.
Uniform boundedness implies a fixed positive fraction of the remaining
gap is removed over each window of length $M$.
\end{proof}

% ------------------------------------------------------------

\begin{corollary}[Density Exhaustion]
\label{cor:density-exhaustion}
Under the hypotheses of Theorem~\ref{thm:gap-contraction},
\[
\lim_{m\to\infty} C(m)=1.
\]
\end{corollary}

\begin{proof}
Iterating the contraction inequality gives
\[
R(m+kM)\le (1-\epsilon)^kR(m)\to0.
\]
Hence $C(m)\to1$.
\end{proof}

% ------------------------------------------------------------

\begin{remark}[Critical Reduction]
The remaining task is purely combinatorial: prove a uniform bound on the
delay and size of compensation packets generated by jump indices in the
admissible division-word grammar.
\end{remark}

% ------------------------------------------------------------
\section{Grammar Growth and A020914 Structure}
\label{sec:8a-grammar}

\begin{itemize}
    \item Role of A020914 as delay/deficit structure
    \item Role of A186009 as first-descent enumeration
    \item Growth of admissible words
    \item Non-uniform branching effects
\end{itemize}

% ------------------------------------------------------------
\section{Block Decomposition Principle}
\label{sec:8a-blocks}

We formalize a covering mechanism based on grouping contributions
from minimal descent words into finite blocks that collectively
exhaust dyadic mass deficits.

% ------------------------------------------------------------
\begin{definition}[Residual Gap]
Let
\[
C(n) = \sum_{k=0}^n N(k),
\]
and define the residual gap
\[
R(n) := 1 - C(n).
\]
\end{definition}

% \begin{definition}[Dyadic Deficit]
% Let
% \[
% C(n) := \sum_{k=0}^n N(k),
% \]
% where $N(k)$ denotes the contribution of minimal descent words at level $k$.

% Define the dyadic deficit at level $k$ by
% \[
% \Delta_k := 2^{-k}.
% \]
% \end{definition}

% ------------------------------------------------------------

\begin{definition}[Block Interval]
A block interval is a contiguous index set
\[
B = \{n_0, n_0+1, \dots, n_0+\ell-1\}
\]
of length $\ell \ge 1$.
\end{definition}

% ------------------------------------------------------------

% \begin{definition}[Finite Recovery Window]
% We say the system admits a finite recovery window if there exists
% an integer $M \ge 1$ such that for every $k$ there exists a block
% \[
% B_k = \{n_k, n_k+1, \dots, n_k+M-1\}
% \]
% satisfying
% \[
% \sum_{n \in B_k} N(n) \;\ge\; 2^{-k}.
% \]
% \end{definition}

% ------------------------------------------------------------

\begin{lemma}[Mass Injection Implies Contraction]
\label{lem:mass-injection}
Suppose there exist constants $M \ge 1$ and $\epsilon > 0$ such that
for all sufficiently large $n$,
\[
\sum_{k=n+1}^{n+M} N(k) \;\ge\; \epsilon \, R(n).
\]

Then
\[
R(n+M) \le (1-\epsilon) R(n).
\]
\end{lemma}

\begin{proof}
We compute:
\[
R(n+M)
=
1 - C(n+M)
=
1 - \bigl(C(n) + (C(n+M)-C(n))\bigr).
\]

Using the hypothesis,
\[
C(n+M)-C(n) \ge \epsilon R(n),
\]
so
\[
R(n+M)
\le
1 - C(n) - \epsilon R(n)
=
R(n) - \epsilon R(n)
=
(1-\epsilon) R(n).
\]
\end{proof}

\begin{definition}[Dyadic Contraction Property]
We say the system satisfies dyadic contraction if there exists
an integer $M \ge 1$ such that for all sufficiently large $n$,
\[
R(n+M) \le \tfrac{1}{2} R(n).
\]
\end{definition}

% ------------------------------------------------------------

% \begin{theorem}[Gap Contraction Implies Density Exhaustion]
% \label{thm:gap-exhaustion}
% If there exist constants $M \ge 1$ and $\epsilon > 0$ such that
% \[
% R(n+M) \le (1-\epsilon) R(n)
% \]
% for all sufficiently large $n$, then
% \[
% \lim_{n \to \infty} C(n) = 1.
% \]
% \end{theorem}

% \begin{proof}
% Iterating the contraction,
% \[
% R(n + kM) \le (1-\epsilon)^k R(n).
% \]

% Since $0 < 1-\epsilon < 1$, the right-hand side tends to zero as $k \to \infty$.

% Thus $R(n) \to 0$, and hence $C(n) \to 1$.
% \end{proof}
\begin{theorem}[Dyadic Contraction Implies Density Exhaustion]
\label{thm:dyadic-exhaustion}
If the dyadic contraction property holds, then
\[
\lim_{n \to \infty} C(n) = 1.
\]
\end{theorem}

\begin{proof}
Iterating,
\[
R(n + kM) \le 2^{-k} R(n).
\]

Thus $R(n) \to 0$, and hence $C(n) \to 1$.
\end{proof}

% ------------------------------------------------------------

\begin{remark}[Critical Reduction]
The problem reduces to establishing dyadic contraction over a finite window,
i.e., the existence of $M$ such that
\[
R(n+M) \le \tfrac{1}{2} R(n).
\]
\end{remark}

% ------------------------------------------------------------
\section{Gap Dynamics and Feedback System}
\label{sec:8a-feedback}

\begin{itemize}
    \item Delayed interaction between numerator and denominator growth
    \item A020914-driven structural lag
    \item Bounded deficit propagation
    \item Stabilisation of block contribution size
\end{itemize}

\begin{lemma}[Uniform Mass Injection over Finite Windows]
\label{lem:mass-injection}
There exist constants $M \ge 1$ and $\epsilon > 0$ such that for all sufficiently large $n$,
\[
\sum_{k=n+1}^{n+M} N(k) \;\ge\; \epsilon \, R(n).
\]
\end{lemma}

\begin{definition}[Finite Contraction Window]
We say the system admits a finite contraction window if there exists
an integer $M \ge 1$ such that for all sufficiently large $n$,
\[
R(n+M) \le \tfrac{1}{2} R(n).
\]
\end{definition}

Candidate proof sketch: 
use A020914 structure to identify blocks of minimal
words that collectively inject mass proportional to the current gap.
\begin{definition}[Unresolved Classes]
Let $\mathcal{R}_n$ be the set of residue classes modulo $2^{d(n)}$
corresponding to integers with no descent prefix up to length $n$.
\end{definition}

\begin{proposition}[Bounded Recovery Window from Jump Structure]
There exists a finite constant $M$ such that for all sufficiently large $n$,
\[
\sum_{k=n+1}^{n+M} N(k) \ge \epsilon R(n)
\]
for some $\epsilon > 0$.
\end{proposition}

\begin{lemma}[Jump-Induced Expansion]
Each increase in $d(n)$ forces the creation of at least one new
minimal descent word within a bounded number of steps.
\end{lemma}

\begin{lemma}[Bounded Compensation Packets]
\label{lem:compensation-packets}
There exist constants $L \ge 1$ and $c > 0$ such that
for every jump index $n$, there exists an interval
\[
I_n \subset \{n+1, \dots, n+L\}
\]
satisfying
\[
\sum_{k \in I_n} N(k) \;\ge\; c \cdot 2^{-d(n)}.
\]
\end{lemma}

same?

\begin{lemma}[Jump-Induced Numerator Growth]
\label{lem:jump-num-growth}
There exist constants $L \ge 1$ and $\alpha > 0$ such that
for all sufficiently large $n$,
\[
\sum_{i=1}^{L} A(n+i)
\;\ge\;
\alpha \cdot 2^{J(n,L)}.
\]
\end{lemma}

\begin{corollary}[Uniform Compensation from Jump Structure]
There exist constants $L \ge 1$ and $c > 0$ such that
for all sufficiently large $n$,
\[
\sum_{i=1}^{L} N(n+i)
\;\ge\;
c \cdot 2^{-d(n)}.
\]
\end{corollary}


\begin{proposition}[Finite Compensation Window]
There exist constants $M \ge 1$ and $\epsilon > 0$ such that
for all sufficiently large $n$,
\[
\sum_{k=n+1}^{n+M} N(k)
\;\ge\;
\epsilon R(n).
\]
\end{proposition}


This is cool because it shows that the system has a built-in feedback mechanism
that prevents the gap from persisting indefinitely. Each time the denominator
jumps, it creates a new opportunity for mass injection that helps to close the gap.

% ------------------------------------------------------------
\subsection{Dyadic Cancellation and Effective Mass Constant}
\label{subsec:dyadic-cancellation}

Let $d(n) = A020914(n)$ and $A(n) = A186009(n)$.
Define
\[
N(n) := \frac{A(n+1)}{2^{d(n)}}.
\]

We analyze the effect of jump accumulation on block sums of $N(n)$.

% ------------------------------------------------------------

\begin{lemma}[Denominator Expansion under Jump Decomposition]
\label{lem:denominator-jump-expansion}
Let
\[
\Delta d(n) := d(n+1) - d(n).
\]
Then for any $i \ge 1$,
\[
d(n+i)
=
d(n) + i + J(n,i),
\]
so
\[
2^{d(n+i)}
=
2^{d(n)} \cdot 2^{i} \cdot 2^{J(n,i)},
\]
where
\[
J(n,i) := \sum_{k=0}^{i-1} \bigl(\Delta d(n+k) - 1\bigr)_+
\]
is the accumulated jump excess.
\end{lemma}

% ------------------------------------------------------------

\begin{lemma}[Cancellation Structure of Dyadic Scaling]
\label{lem:dyadic-cancellation}
Let $N(n) = A(n+1)/2^{d(n)}$. Then
\[
N(n+i)
=
\frac{A(n+i+1)}{2^{d(n)} \cdot 2^{i} \cdot 2^{J(n,i)}}.
\]
\end{lemma}

\begin{proof}
Immediate from Lemma~\ref{lem:denominator-jump-expansion}.
\end{proof}

% ------------------------------------------------------------

\begin{lemma}[Jump-Compensated Numerator Growth]
\label{lem:jump-comp-growth}
There exist constants $L \ge 1$ and $\alpha > 0$ such that
for all sufficiently large $n$,
\[
\sum_{i=1}^{L} A(n+i)
\;\ge\;
\alpha \cdot 2^{J(n,L)}.
\]
\end{lemma}

% ------------------------------------------------------------

\begin{proposition}[Dyadic Cancellation Lower Bound]
\label{prop:dyadic-cancellation-bound}
There exist constants $L \ge 1$ and $c > 0$ such that
for all sufficiently large $n$,
\[
\sum_{i=1}^{L} N(n+i)
\;\ge\;
\frac{c}{2^{d(n)}}.
\]
Moreover, the constant $c$ admits the estimate
\[
c
\;\asymp\;
\alpha \sum_{i=1}^{L} 2^{-i},
\]
up to bounded distortion from jump fluctuations.
\end{proposition}

\begin{proof}
Using Lemma~\ref{lem:dyadic-cancellation},
\[
\sum_{i=1}^{L} N(n+i)
=
\sum_{i=1}^{L}
\frac{A(n+i+1)}{2^{d(n)} \cdot 2^{i} \cdot 2^{J(n,i)}}.
\]

Applying Lemma~\ref{lem:jump-comp-growth} yields
\[
A(n+i+1) \gtrsim 2^{J(n,i)},
\]
so the jump terms cancel at leading order:
\[
\sum_{i=1}^{L} N(n+i)
\gtrsim
\frac{1}{2^{d(n)}} \sum_{i=1}^{L} 2^{-i}.
\]

Hence
\[
c \sim \alpha \sum_{i=1}^{L} 2^{-i},
\]
with bounded correction from non-uniform jump distribution.
\end{proof}


More stuff...

% ============================================================
\begin{lemma}[A020914 Jump Measure Decomposition]
\label{lem:a020914-measure-decomp}
% ============================================================

Let $d(n)=A020914(n)$ and define the dyadic cylinder at level $n$ by
\[
\mathcal{I}_n := \{ \text{admissible division words with prefix length } n \}.
\]

Let $\mu_n$ denote the induced mass assignment
\[
\mu_n := \sum_{w \in \mathcal{I}_n} N(w),
\quad \text{so that } C(n)=\sum_{k\le n} N(k).
\]

Suppose $n$ is a jump index, i.e.
\[
d(n+1)=d(n)+2.
\]

Then the refinement from level $n$ to level $n+1$ induces a
strict partition refinement of dyadic cylinders:
\[
\mathcal{I}_n \;\longrightarrow\; \mathcal{I}_{n+1}^{(1)} \cup \mathcal{I}_{n+1}^{(2)},
\]
where:

\begin{enumerate}
\item $\mathcal{I}_{n+1}^{(1)}$ consists of cylinders already present under unit-step growth (the “stable extension”),
\item $\mathcal{I}_{n+1}^{(2)}$ consists of newly created admissible cylinders induced by the jump.
\end{enumerate}

Moreover, this decomposition is disjoint:
\[
\mathcal{I}_{n+1}^{(1)} \cap \mathcal{I}_{n+1}^{(2)} = \varnothing,
\quad
\mathcal{I}_{n+1}^{(1)} \cup \mathcal{I}_{n+1}^{(2)} = \mathcal{I}_{n+1}.
\]

Finally, the jump component satisfies the scaling relation
\[
\mu\big(\mathcal{I}_{n+1}^{(2)}\big)
=
\theta_n \, \mu\big(\mathcal{I}_n^{\mathrm{unresolved}}\big),
\]
for some $\theta_n \in (0,1]$, where
\[
\mathcal{I}_n^{\mathrm{unresolved}} := \mathcal{I}_{\infty} \setminus \bigcup_{k\le n} \mathcal{I}_k.
\]

\end{lemma}

% ============================================================
\section{Grammar–Drift–Block Framework}
\label{sec:grammar-drift-block}
% ============================================================

\begin{table}[htbp]
\centering
\caption{Minimum block size, $M$, for logarithmic convexity of $a(n) = A186009(n)$}
\label{tab:A186009}
\vspace{4pt}
\begin{tabular}{@{}rccrr@{}}
\toprule
$M$ & $a_{n} a_{n+2M} \ge a_{n+M}$ & {S/D} & $n$ & $a(n)$ \\
\midrule
1 & $a_{ 1} a_{ 3} \ge a_{ 2} a_{ 2}$ & S &  1 &            1 \\
1 & $a_{ 2} a_{ 4} \ge a_{ 3} a_{ 3}$ & S &  2 &            1 \\
2 & $a_{ 3} a_{ 7} \ge a_{ 5} a_{ 5}$ & S &  3 &            1 \\
1 & $a_{ 4} a_{ 6} \ge a_{ 5} a_{ 5}$ & D &  4 &            2 \\
2 & $a_{ 5} a_{ 9} \ge a_{ 7} a_{ 7}$ & S &  5 &            3 \\
1 & $a_{ 6} a_{ 8} \ge a_{ 7} a_{ 7}$ & D &  6 &            7 \\
1 & $a_{ 7} a_{ 9} \ge a_{ 8} a_{ 8}$ & S &  7 &           12 \\
3 & $a_{ 8} a_{14} \ge a_{11} a_{11}$ & D &  8 &           30 \\
1 & $a_{ 9} a_{11} \ge a_{10} a_{10}$ & D &  9 &           85 \\
2 & $a_{10} a_{14} \ge a_{12} a_{12}$ & S & 10 &          173 \\
1 & $a_{11} a_{13} \ge a_{12} a_{12}$ & D & 11 &          476 \\
1 & $a_{12} a_{14} \ge a_{13} a_{13}$ & S & 12 &          961 \\
4 & $a_{13} a_{21} \ge a_{17} a_{17}$ & D & 13 &         2652 \\
1 & $a_{14} a_{16} \ge a_{15} a_{15}$ & D & 14 &         8045 \\
3 & $a_{15} a_{21} \ge a_{18} a_{18}$ & S & 15 &        17637 \\
1 & $a_{16} a_{18} \ge a_{17} a_{17}$ & D & 16 &        51033 \\
2 & $a_{17} a_{21} \ge a_{19} a_{19}$ & S & 17 &       108950 \\
1 & $a_{18} a_{20} \ge a_{19} a_{19}$ & D & 18 &       312455 \\
1 & $a_{19} a_{21} \ge a_{20} a_{20}$ & S & 19 &       663535 \\
4 & $a_{20} a_{28} \ge a_{24} a_{24}$ & D & 20 &      1900470 \\
1 & $a_{21} a_{23} \ge a_{22} a_{22}$ & D & 21 &      5936673 \\
2 & $a_{22} a_{26} \ge a_{24} a_{24}$ & S & 22 &     13472296 \\
1 & $a_{23} a_{25} \ge a_{24} a_{24}$ & D & 23 &     39993895 \\
1 & $a_{24} a_{26} \ge a_{25} a_{25}$ & S & 24 &     87986917 \\
4 & $a_{25} a_{33} \ge a_{29} a_{29}$ & D & 25 &    257978502 \\
1 & $a_{26} a_{28} \ge a_{27} a_{27}$ & D & 26 &    820236724 \\
3 & $a_{27} a_{33} \ge a_{31} a_{31}$ & S & 27 &   1899474678 \\
1 & $a_{28} a_{30} \ge a_{29} a_{29}$ & D & 28 &   5723030586 \\
2 & $a_{29} a_{33} \ge a_{31} a_{31}$ & S & 29 &  12809477536 \\
1 & $a_{30} a_{32} \ge a_{31} a_{31}$ & D & 30 &  38036848410 \\
1 & $a_{31} a_{33} \ge a_{32} a_{32}$ & S & 31 &  84141805077 \\
4 & $a_{32} a_{40} \ge a_{36} a_{36}$ & D & 32 & 248369601964 \\
1 & $a_{33} a_{35} \ge a_{34} a_{34}$ & D & 33 & 794919136728 \\
\bottomrule
\end{tabular}
\end{table}

This section formalises the interaction between the combinatorial
structure induced by $A020914(n)$ and the induced growth dynamics
of $C(n) = \sum_{k=0}^n N(k)$.

The central idea is that admissible division words generate a
finite-state grammar, which induces a bounded drift automaton.
This in turn yields a uniform block inequality governing residual decay.

% ------------------------------------------------------------
\subsection{Grammar Layer}
% ------------------------------------------------------------

\begin{definition}[Event Alphabet]
Each index $n$ is assigned an event type:
\[
E(n) \in \{S, D\},
\]
where:
\begin{itemize}
\item $S$ denotes a \emph{single step} (unit increase in $d(n)$),
\item $D$ denotes a \emph{double step} (jump increment in $d(n)$).
\end{itemize}
\end{definition}

\begin{lemma}[Jump-Bounded Grammar]
\label{lem:grammar-bound}
There exists an integer $L \ge 1$ such that no admissible sequence
contains more than $L$ consecutive $D$ events.
\end{lemma}

\begin{remark}
Empirically $L \in \{2,3\}$ dominates all observed regimes, with rare
local fluctuations, but boundedness is the structural requirement.
\end{remark}

% ------------------------------------------------------------
\subsection{Drift Automaton Layer}
% ------------------------------------------------------------

\begin{definition}[Drift Assignment]
Associate to each event a drift contribution:
\[
\Delta(S) = +\alpha,\qquad \Delta(D) = -\beta,
\]
where $\alpha,\beta > 0$ encode numerator and denominator scaling
effects respectively.
\end{definition}

\begin{definition}[Cumulative Drift]
For a block $B = \{n,\dots,n+m\}$ define
\[
\mathcal{D}(B) := \sum_{i=n}^{n+m} \Delta(E(i)).
\]
\end{definition}

\begin{lemma}[Bounded Negative Streaks]
\label{lem:bounded-streaks}
By Lemma~\ref{lem:grammar-bound}, every block of length $L+1$
contains at least one $S$ event.
\end{lemma}

\begin{lemma}[Local Drift Compensation]
\label{lem:local-comp}
If
\[
\alpha \ge \frac{L}{L+1}\beta,
\]
then for every block of length $L+1$,
\[
\mathcal{D}(B) \ge 0.
\]
\end{lemma}

\begin{remark}
This is the first point where grammar constraints convert directly
into quantitative positivity.
\end{remark}

% ------------------------------------------------------------
\subsection{Block Inequality Layer}
% ------------------------------------------------------------

% \begin{lemma}[Block Non-Negativity]
% \label{lem:block-nonneg}
% There exists $M = L+1$ such that for all sufficiently large $n$,
% \[
% C(n+M) - C(n) \ge 0.
% \]
% \end{lemma}

% \begin{proof}
% Each block of length $M$ contains at least one $S$ event and at most
% $L$ $D$ events. Under Lemma~\ref{lem:local-comp}, the net drift is
% non-negative. Since $C(n)$ accumulates drift contributions,
% the inequality follows.
% \end{proof}

\begin{lemma}[Fractional Residual Depletion]
\label{lem:fractional-depletion}
Suppose there exists a block size $M$ and a constant $\eta > 0$ such that
each block introduces a family of disjoint admissible cylinders whose
total measure satisfies
\[
\sum_{n \in B_k} N(n) \ge \eta \, R(n_k),
\]
where $R(n) = 1 - C(n)$.
Then
\[
R(n_k + M) \le (1 - \eta)\,R(n_k).
\]
\end{lemma}

\begin{lemma}[Residual Contraction]
\label{lem:residual-contraction}
Define $R(n) = 1 - C(n)$. Then
\[
R(n+M) \le R(n).
\]
Moreover, if strict positivity occurs in infinitely many blocks,
then $R(n)$ is strictly decreasing on a subsequence.
\end{lemma}


\begin{theorem}[Exact Double-Jump Window Law]
Let $s_k$ denote the start indices of consecutive jump pairs
$(2,2)$ in the increment word $d(n+1)-d(n)$.

Suppose admissible grammar cycles have minimum return length $5$
and maximum return length $7$.

Then every interval of length $8$ contains exactly one $s_k$.
\end{theorem}

\begin{theorem}[Anchored Eight-Step Recovery]
Let $s_k$ be the initial index of a consecutive jump pair $(2,2)$
in the increment word $d(n+1)-d(n)$.

Then the local convexity defect generated at $s_k$
is repaired within eight steps, in the sense that
\[
a_{s_k}a_{s_k+8}\ge a_{s_k+4}^2.
\]
\end{theorem}

% ------------------------------------------------------------
\subsection{Main Structural Theorem}
% ------------------------------------------------------------

\begin{theorem}[Grammar-Induced Block Exhaustion]
\label{thm:block-exhaustion-final}
If:
\begin{enumerate}
\item the event grammar is jump-bounded (Lemma~\ref{lem:grammar-bound}),
\item drift contributions satisfy $\alpha > 0$, $\beta > 0$,
\item and local compensation holds (Lemma~\ref{lem:local-comp}),
\end{enumerate}
then there exists $M < \infty$ such that
\[
\lim_{n \to \infty} C(n) = 1.
\]
\end{theorem}

\begin{proof}
By Lemma~\ref{lem:block-nonneg}, the system admits non-negative drift
on every block of length $M$. Hence the residual $R(n)$ is monotone
non-increasing and bounded below by $0$.

Since each block contains at least one compensating event with
strict contribution in infinitely many positions, the residual cannot
stabilise above a positive constant. Therefore $R(n) \to 0$.
\end{proof}

% ------------------------------------------------------------
\section{Grammar-Driven Recovery Dynamics of the Dyadic Threshold Sequence}
% ------------------------------------------------------------

\begin{definition}[Dyadic Threshold Sequence]
Let
\[
d(n)=A020914(n)=\lfloor n\log_2 3\rfloor+1.
\]
Define the increment word
\[
\delta_n=d(n+1)-d(n)\in\{1,2\}.
\]
A symbol $\delta_n=2$ is called a jump.
\end{definition}

\begin{definition}[Double-Jump Event]
A double-jump event occurs at index $n$ if
\[
\delta_n=\delta_{n+1}=2.
\]
Its start index is denoted $s_k$.
\end{definition}

\begin{lemma}[Minimum Return Length]
For consecutive start indices,
\[
s_{k+1}-s_k\ge5.
\]
\end{lemma}

\begin{lemma}[Maximum Return Length]
For consecutive start indices,
\[
s_{k+1}-s_k\le7.
\]
\end{lemma}

\begin{corollary}[Shock-Free Midpoint Window]
For every $k$, no new double-jump event begins in
\[
\{s_k+1,\dots,s_k+4\}.
\]
\end{corollary}

\begin{corollary}[Forced Recurrence Window]
For every $k$, at least one double-jump event begins in
\[
\{s_k+5,\dots,s_k+7\}.
\]
\end{corollary}

\begin{definition}[Recovery Sequence]
Let
\[
a(n)=A186009(n),
\]
the number of new first-descent words at level $n$.
\end{definition}

\begin{lemma}[Grammar Compensation Principle]
Each jump pair $(2,2)$ initially suppresses local dyadic mass,
but the forced recurrence window creates additional admissible
first-descent words whose aggregate contribution compensates
the earlier deficit within at most eight steps.
\end{lemma}

\begin{theorem}[Anchored Eight-Step Recovery]
For every double-jump start index $s_k$,
\[
a(s_k)\,a(s_k+8)\ge a(s_k+4)^2.
\]
\end{theorem}

\begin{corollary}[Uniform Block Convexity]
There exists $M\le4$ such that for all sufficiently large $n$,
\[
a(n)a(n+2M)\ge a(n+M)^2.
\]
\end{corollary}

\begin{theorem}[Dyadic Gap Contraction]
The grammar-driven recovery mechanism implies the residual gap
contracts on bounded windows:
\[
R(n+M)\le \frac12 R(n).
\]
Hence
\[
\lim_{n\to\infty}C(n)=1.
\]
\end{theorem}


\begin{theorem}[Bounded Dyadic Shell Capture]
There exists an integer $L\ge1$ such that for every sufficiently large $n$,
if
\[
2^{-(m+1)} < R(n)\le 2^{-m},
\]
then
\[
\sum_{j=n}^{n+L} N(j)\ge 2^{-(m+1)}.
\]
Consequently,
\[
R(n+L)\le \frac12 R(n).
\]
\end{theorem}


\begin{theorem}[Uniform Twelve-Block Recurrence]
Every interval of twelve consecutive indices contains at least two
double-jump starts.
\end{theorem}


\subsection{Cycle Structure and Net Multiplicative Drift}

Let
\[
\delta_n=d(n+1)-d(n)\in\{1,2\},
\]
where $d(n)=A020914(n)$.
We write
\[
s:=1,\qquad d:=2,
\]
for single and double dyadic increments.

For a block word
\[
w=(\delta_n,\dots,\delta_{n+L-1}),
\]
define its total dyadic increment
\[
\Sigma(w)=\sum_{j=0}^{L-1}\delta_{n+j},
\]
and associated drift ratio
\[
\mu(w)=\frac{3^L}{2^{\Sigma(w)}}.
\]

The block is called contractive if $\mu(w)<1$, and expansive if $\mu(w)>1$.

A frequently occurring grammar block is the $7$-cycle
\[
(s,d,s,d,s,d,d),
\]
for which
\[
\Sigma=1+2+1+2+1+2+2=11.
\]
Hence
\[
\mu_7=\frac{3^7}{2^{11}}<1,
\]
so the $7$-cycle is contractive.

Another admissible block is the $5$-cycle
\[
(s,d,s,d,d),
\]
with
\[
\Sigma=1+2+1+2+2=8.
\]
Thus
\[
\mu_5=\frac{3^5}{2^8}>1,
\]
so the $5$-cycle is locally expansive.

Concatenating these gives the forced $12$-cycle
\[
(s,d,s,d,s,d,d,s,d,s,d,d),
\]
with total increment
\[
\Sigma=11+8=19.
\]
Therefore
\[
\mu_{12}=\frac{3^{12}}{2^{19}}<1.
\]

Hence the grammar-induced $12$-cycle has negative net drift:
although the $5$-cycle is locally expansive, it is dominated by the
preceding $7$-cycle. Repeated admissible blocks therefore exhibit
average contraction.

% ------------------------------------------------------------
\section{Denominator Beattic Composition and Jump-Compensated Numerator Growth}

\begin{lemma}[Diminishing Tower Error]
Let $\alpha=\log_2 3$, and let $q_k$ be the ordered sequence of convergent and semiconvergent denominators associated to $\alpha$.

Then
\[
\epsilon_k := q_k\alpha - p_k
\]
for nearest integers $p_k$ satisfies
\[
|\epsilon_k|\to 0.
\]

Hence the corresponding tower blocks become asymptotically exact dyadic/ternary synchronizations.
\end{lemma}


% ------------------------------------------------------------
\section{Beatty Structure and Tower Decomposition of the Denominator}
\label{sec:beatty-towers}

We formalize the macroscopic structure governing the denominator growth
through the Beatty sequence associated with $\log_2 3$.

% ------------------------------------------------------------

\begin{definition}[Dyadic Threshold Sequence]
Let
\[
\alpha := \log_2 3,
\qquad
d(n) := A020914(n) = \lfloor n\alpha \rfloor + 1.
\]
\end{definition}

The increment word
\[
\Delta(n) := d(n+1)-d(n)
\]
takes values in $\{1,2\}$ and defines a Sturmian (Beatty) sequence.

% ------------------------------------------------------------

\begin{definition}[Tower Lengths]
A sequence of integers $\{q_k\}$ is called a tower sequence if
\[
q_k \alpha = p_k + \varepsilon_k,
\]
where $p_k \in \mathbb{Z}$ and $|\varepsilon_k|$ is minimal among all
denominators up to $q_k$.

We refer to $q_k$ as a \emph{tower length} and $p_k$ as its associated
dyadic exponent.
\end{definition}

Empirically, the tower sequence begins as
\[
5,\,7,\,12,\,41,\,53,\,306,\,359,\,665,\,15601,\,16266,\,31867,\,\dots
\]

% ------------------------------------------------------------

\begin{lemma}[Diminishing Approximation Error]
\label{lem:tower-error}
Let $\varepsilon_k := q_k \alpha - p_k$.
Then
\[
|\varepsilon_k| \to 0 \quad \text{as } k \to \infty.
\]
\end{lemma}

\begin{proof}
This follows from the theory of continued fractions and best rational
approximations of irrational numbers.
\end{proof}

% ------------------------------------------------------------

\begin{lemma}[Alternating Approximation Sign]
\label{lem:alternating}
The errors $\varepsilon_k$ alternate in sign.  That is,
tower approximations alternate between underestimates and overestimates
of $\alpha$.
\end{lemma}

\begin{proof}
This is a standard property of convergents of continued fractions.
\end{proof}

% ------------------------------------------------------------

\begin{definition}[Expansive and Contractive Towers]
A tower length $q_k$ is called:
\begin{itemize}
\item \emph{expansive} if $\varepsilon_k < 0$ (i.e.\ $q_k\alpha < p_k$),
\item \emph{contractive} if $\varepsilon_k > 0$.
\end{itemize}
\end{definition}

Thus expansive towers slightly undercount dyadic growth, while
contractive towers slightly overcount it.

% ------------------------------------------------------------

\begin{lemma}[Tower Partition Structure]
\label{lem:tower-partition}
Tower lengths occur in finite clusters of the form
\[
\{5,7,12\},\quad
\{41,53\},\quad
\{306,359,665\},\quad
\{15601,16266,31867\},\ \dots
\]
in which:
\begin{itemize}
\item the smallest element is expansive,
\item all remaining elements are contractive.
\end{itemize}
\end{lemma}

\begin{proof}[Heuristic justification]
These clusters arise from convergents and semiconvergents of $\alpha$.
The smallest element in each cluster lies below $\alpha$ (negative error),
while subsequent refinements overshoot (positive error) before the next
cycle begins.
\end{proof}

% ------------------------------------------------------------

\begin{lemma}[Exact Jump Count on Towers]
\label{lem:exact-jumps}
Let $q_k$ be a tower length.  Then the increment word over the interval
$[n,n+q_k)$ contains an exact number of double steps (jumps), independent
of $n$.
\end{lemma}

\begin{proof}[Sketch]
Since
\[
d(n+q_k)-d(n) \in \{p_k, p_k+1\},
\]
and $\varepsilon_k$ is small, the total number of increments equal to $2$
is fixed across the block.  The Sturmian structure ensures that all
length-$q_k$ factors have identical Parikh counts.
\end{proof}

% ------------------------------------------------------------

\begin{lemma}[Terminal Double--Double Structure]
\label{lem:dd-terminal}
Every tower block ends in a double--double pattern $(2,2)$.
\end{lemma}

\begin{proof}[Heuristic]
The atomic building blocks of the increment word are the $5$-cycle and
$7$-cycle, both of which terminate in $(2,2)$.  Since all tower lengths
are concatenations of these atomic cycles, the terminal structure is
preserved.
\end{proof}

% ------------------------------------------------------------

\begin{theorem}[Macroscopic Denominator Law]
\label{thm:denominator-law}
Let $q_k$ be a tower sequence.  Then over any interval of length $q_k$:
\begin{enumerate}
\item The total dyadic exponent satisfies
\[
d(n+q_k)-d(n) = p_k + O(1),
\]
with $|O(1)| \le 1$.

\item The jump count is exact and independent of $n$.

\item The multiplicative imbalance satisfies
\[
\frac{3^{q_k}}{2^{p_k}} = 2^{\varepsilon_k},
\]
where $\varepsilon_k \to 0$.

\item Expansive towers ($\varepsilon_k<0$) create slack, while
contractive towers ($\varepsilon_k>0$) consume slack.

\item The sequence of towers alternates between expansion and contraction,
with diminishing magnitude.
\end{enumerate}
\end{theorem}

% ------------------------------------------------------------

\begin{remark}[Denominator Rigidity]
The denominator is completely governed by the Beatty structure of
$\log_2 3$.  All macroscopic irregularities reduce to a controlled,
alternating sequence of small multiplicative errors across tower scales.
\end{remark}

% ------------------------------------------------------------
\section{Renomalization Invariance and Self-Replication of the Admissible Division Word Distribution}
\label{sec:ren-invariance}

\begin{theorem}[Exact Renormalization Invariance of $A186009$]
Let $v_n$ denote the tuple of admissible division word counts
contributing to $A186009(n)$.

Then there exists a normalized vector $v$ such that:

\[
v = \frac{T(v)}{\|T(v)\|_1},
\]

and for all $n$,
\[
v_n = \|v_n\|_1 \cdot v.
\]

In particular, the sequence is self-replicating under normalization.
\end{theorem}

\subsection{Beatty Structure and Uniform Jump Control}
\label{sec:beatty-structure}

The dyadic threshold sequence
\[
d(n) = A020914(n)
\]
admits the explicit representation
\[
d(n) = \lfloor n \log_2 3 \rfloor + 1.
\]
Thus the increment sequence
\[
\Delta d(n) := d(n+1) - d(n)
\]
takes values in $\{1,2\}$ and forms a Sturmian word of slope
\[
\alpha = \log_2 3.
\]

\begin{lemma}[Sturmian Balance]
\label{lem:sturmian-balance}
For every interval of length $L$,
\[
\sum_{k=0}^{L-1} \bigl(\Delta d(n+k) - \alpha\bigr)
= O(1),
\]
where the implied constant is independent of $n$ and $L$.
\end{lemma}

\begin{proof}
This is a standard property of Beatty (Sturmian) sequences arising from
irrational rotations. The discrepancy between the discrete sum and its
linear expectation is uniformly bounded.
\end{proof}

\begin{corollary}[Uniform Dyadic Growth]
\label{cor:uniform-dyadic}
For all $n$ and $L$,
\[
d(n+L) - d(n) = L \log_2 3 + O(1).
\]
Equivalently,
\[
2^{d(n+L)} = 2^{d(n)} \cdot 3^L \cdot \Theta(1),
\]
where the implicit constants are absolute.
\end{corollary}

\subsection{Tower Indices and Exact Calibration}
\label{sec:tower-indices}

Let $\{q_k\}$ denote the sequence of convergents to $\log_2 3$ arising from
its continued fraction expansion. At these indices,
\[
q_k \log_2 3 = p_k + \varepsilon_k,
\quad |\varepsilon_k| \to 0.
\]

\begin{lemma}[Exact Jump Count at Tower Indices]
\label{lem:tower-exact}
At each tower index $q_k$,
\[
d(q_k) = p_k,
\]
up to a bounded additive constant, and the number of jump increments
$\Delta d(n)=2$ up to $q_k$ is determined exactly.
\end{lemma}

\begin{remark}
Tower indices provide exact realizations of the dyadic–ternary balance.
Lemma~\ref{lem:sturmian-balance} shows that all other indices deviate from
this ideal behavior by at most a uniformly bounded error.
\end{remark}

\subsection{Transfer to the Numerator}
\label{sec:numerator-transfer}

Let
\[
a(n) = A186009(n),
\qquad
N(n) = \frac{a(n)}{2^{d(n)}}.
\]

The admissible grammar induces new division words precisely at jump events,
but with a bounded delay.

\begin{lemma}[Delayed Sturmian Inheritance]
\label{lem:delay-inheritance}
There exists $L_0 \ge 1$ such that the growth of $a(n)$ satisfies
\[
a(n+L) = a(n)\cdot 3^L \cdot \Theta(1)
\]
for all $L \ge L_0$, uniformly in $n$.
\end{lemma}

\begin{proof}[Sketch]
New admissible division words are created by expansions of the prefix-sum
constraint, which occur at jump indices of $d(n)$. Since these jumps follow
a Sturmian distribution and every jump produces new words within bounded
delay, the cumulative growth of $a(n)$ tracks the same exponential rate
$3^n$ up to multiplicative constants.
\end{proof}

\subsection{Block Convexity and Exponential Envelope}
\label{sec:block-convexity}

\begin{lemma}[Block Log-Convexity]
\label{lem:block-convex}
There exists $M \ge 1$ such that for all $n$,
\[
a(n)\,a(n+2M) \ge a(n+M)^2.
\]
\end{lemma}

\begin{proof}[Sketch]
This follows from the bounded delay in numerator response together with
the Sturmian regularity of jump distribution. Local fluctuations are
smoothed over blocks of length $M$, yielding convexity at that scale.
\end{proof}

\begin{corollary}[Exponential Envelope]
\label{cor:exp-envelope}
There exists $\lambda > 0$ such that for all $n$ and $k \ge 0$,
\[
a(n+k) = a(n)\,\lambda^k \cdot \Theta(1).
\]
\end{corollary}

\subsection{Measure Decomposition and Geometric Decay}
\label{sec:measure-decay}

Combining Corollaries~\ref{cor:uniform-dyadic} and
\ref{cor:exp-envelope}, we obtain
\[
N(n+k)
= \frac{a(n+k)}{2^{d(n+k)}}
= \frac{a(n)}{2^{d(n)}} \left(\frac{\lambda}{3}\right)^k \cdot \Theta(1).
\]

\begin{lemma}[Block Ratio Law]
\label{lem:block-ratio}
There exists $L \ge 1$ such that for all $n$,
\[
\frac{\sum_{k=L}^{2L-1} N(n+k)}{\sum_{k=0}^{L-1} N(n+k)}
= \rho + o(1),
\]
where $0 < \rho < 1$ is a constant.
\end{lemma}

\begin{remark}
Empirically, $\rho \approx \tfrac{1}{2}$, but only the strict inequality
$\rho < 1$ is required for the arguments below.
\end{remark}

\begin{theorem}[Uniform Mass Contraction]
\label{thm:mass-contraction}
There exist constants $M \ge 1$ and $\epsilon > 0$ such that
\[
R(n+M) \le (1-\epsilon)\,R(n)
\]
for all sufficiently large $n$, where
\[
R(n) = 1 - \sum_{k=0}^{n} N(k).
\]
\end{theorem}

\begin{proof}
By Lemma~\ref{lem:block-ratio}, each block of length $M$ captures a fixed
fraction of the remaining mass. Iterating yields geometric decay of $R(n)$.
\end{proof}

\begin{theorem}[Density Exhaustion]
\label{thm:density-exhaustion-final}
\[
\lim_{n \to \infty} \sum_{k=0}^{n} N(k) = 1.
\]
\end{theorem}

% ============================================================
\section{Canonical Minimal Envelope and Regenerative Growth}
\label{sec:minimal-envelope}
% ============================================================

Let
\[
a(n)=A186009(n),
\]
and let
\[
d(n)=A020914(n)=\lfloor n\log_2 3\rfloor+1.
\]

Define the increment sequence
\[
\delta(n):=d(n+1)-d(n)\in\{1,2\}.
\]

A value $\delta(n)=2$ induces a terminal duplication event in the
admissible tuple grammar.

The central structural observation is that the admissible grammar is
monotone with respect to the placement of duplication events:
earlier duplication permanently increases future tuple mass.

Consequently, the canonical Beatty phase beginning at intercept $0$
produces the minimal admissible growth envelope.

% ------------------------------------------------------------

\subsection{Jump Schedules and Tuple Evolution}

% ------------------------------------------------------------

\begin{definition}[Jump Schedule]
A jump schedule is a sequence
\[
J=(\delta_0,\delta_1,\delta_2,\dots),
\qquad
\delta_i\in\{1,2\},
\]
describing the placement of terminal duplication events in the
tuple recursion.

The canonical schedule is the Beatty/Sturmian increment sequence
\[
\delta_i=d(i+1)-d(i),
\]
associated to
\[
d(i)=\lfloor i\log_2 3\rfloor+1.
\]
\end{definition}

% ------------------------------------------------------------

\begin{definition}[Generating Tuple]
Let
\[
T_n=(t_1,t_2,\dots,t_r)
\]
denote the admissible generating tuple at level $n$.

The tuple evolves by:
\begin{enumerate}
\item cumulative Pascal-type addition,
\item terminal duplication whenever $\delta(n)=2$.
\end{enumerate}

Define the total tuple mass by
\[
|T_n|:=\sum_i t_i.
\]
\end{definition}

% ------------------------------------------------------------

\subsection{Monotonicity of the Grammar}

% ------------------------------------------------------------

\begin{lemma}[Positivity Preservation]
\label{lem:positivity}
All tuple entries remain positive under admissible evolution.

Consequently, any increase introduced at one stage propagates
monotonically to all future stages.
\end{lemma}

\begin{proof}
The admissible recursion consists entirely of cumulative additions
of positive quantities.

Terminal duplication copies an existing positive terminal entry,
and no subtraction occurs anywhere in the grammar evolution.

Hence all tuple entries remain positive, and any increase introduced
at one stage propagates to all future descendants.
\end{proof}

% ------------------------------------------------------------

\begin{lemma}[Earlier Duplication Dominance]
\label{lem:earlier-duplication}
Let
\[
J_1,\qquad J_2
\]
be two admissible jump schedules that agree through some index $m$.

Suppose $J_1$ performs a terminal duplication strictly earlier than
$J_2$.

Then the corresponding tuple masses satisfy
\[
|T_n^{(1)}|
\ge
|T_n^{(2)}|
\]
for all sufficiently large $n$, with strict inequality eventually
persistent.
\end{lemma}

\begin{proof}[Heuristic mechanism]
Earlier duplication inserts additional positive mass into the tuple
recursion sooner.

By Lemma~\ref{lem:positivity}, the admissible evolution is monotone:
all future cumulative sums inherit any earlier increase.

Thus earlier duplication permanently amplifies future tuple growth,
yielding eventual domination.
\end{proof}

% ------------------------------------------------------------

\subsection{Canonical Minimal Envelope}

% ------------------------------------------------------------

\begin{theorem}[Canonical Minimal Envelope]
\label{thm:minimal-envelope}
Among all admissible phase continuations of the Sturmian jump grammar,
the canonical intercept-$0$ Beatty phase produces minimal tuple growth.

Equivalently, for every admissible phase shift $\phi$,
\[
|T_n^{(0)}|
\le
|T_n^{(\phi)}|
\]
for all sufficiently large $n$.
\end{theorem}

\begin{proof}[Structural mechanism]
The canonical Beatty phase delays terminal duplication events
as long as admissibly possible.

By Lemma~\ref{lem:earlier-duplication}, any admissible phase shift
that introduces duplication earlier produces permanently larger
future tuple mass.

Hence the canonical phase yields the minimal admissible growth
envelope.
\end{proof}

% ------------------------------------------------------------

\begin{remark}[Minimal Regenerative Orbit]
Exact self-replication of admissible tuples occurs only after
renormalization together with resetting the Beatty phase to the
canonical intercept-$0$ schedule.

However, Theorem~\ref{thm:minimal-envelope} shows that this reset
produces the smallest admissible continuation.

Thus exact regeneration is unnecessary:
the canonical reset already provides a universal lower envelope for
future admissible growth.
\end{remark}

% ------------------------------------------------------------

\subsection{Regenerative Lower Bound}

% ------------------------------------------------------------

\begin{theorem}[Regenerative Supermultiplicativity]
\label{thm:supermultiplicative}
For all admissible indices $m,n$,
\[
a(n+m)\ge a(n)a(m).
\]
\end{theorem}

\begin{proof}[Heuristic mechanism]
Consider a generating tuple contributing total mass
\[
a(n)=|T_n|.
\]

Discard the internal tuple structure and retain only the aggregate
mass.

After normalization, restart the admissible grammar using the
canonical intercept-$0$ Beatty phase.

By Theorem~\ref{thm:minimal-envelope}, this regenerated evolution
produces the smallest admissible future growth among all admissible
continuations.

But the canonical regeneration reproduces exactly the admissible
growth sequence
\[
a(m).
\]

Hence every admissible configuration contributing to $a(n)$
generates at least $a(m)$ admissible descendants after $m$
additional steps.

Therefore,
\[
a(n+m)\ge a(n)a(m).
\]
\end{proof}

% ------------------------------------------------------------

\subsection{Existence of an Exponential Scale}

% ------------------------------------------------------------

\begin{theorem}[Existence of the Exponential Growth Constant]
\label{thm:lambda-exists}
There exists a finite constant
\[
\lambda>0
\]
such that
\[
\lambda
=
\lim_{n\to\infty} a(n)^{1/n}.
\]
Equivalently,
\[
\lim_{n\to\infty}\frac{\log a(n)}{n}
=
\log\lambda.
\]
\end{theorem}

\begin{proof}
Define
\[
b(n):=\log a(n).
\]

By Theorem~\ref{thm:supermultiplicative},
\[
b(n+m)\ge b(n)+b(m),
\]
so the sequence $b(n)$ is superadditive.

Therefore Fekete's lemma applies, yielding existence of the limit
\[
\lim_{n\to\infty}\frac{b(n)}{n}
=
\sup_n \frac{b(n)}{n}.
\]

Exponentiating gives
\[
\lim_{n\to\infty} a(n)^{1/n}
=
\lambda
\]
for some finite constant $\lambda>0$.
\end{proof}

% ------------------------------------------------------------

\begin{remark}[Delayed Compensation and Convexity]
The canonical Beatty phase minimizes growth by delaying terminal
duplication events as long as admissibly possible.

Consequently, local suppression of tuple growth necessarily forces
larger future compensation once duplication eventually occurs.

This delayed-compensation mechanism naturally induces:
\begin{itemize}
\item block log-convexity,
\item bounded recovery windows,
\item asymptotic smoothing of local oscillations.
\end{itemize}
\end{remark}

% ------------------------------------------------------------

\begin{remark}[Structural Interpretation]
The admissible division-word grammar defines a positive monotone
dynamical system controlled by a Sturmian jump schedule.

The canonical intercept-$0$ Beatty phase acts as the unique minimal
growth orbit, while all admissible phase shifts dominate it.

Thus the existence of the exponential scale $\lambda$ emerges from
an extremal monotone regeneration principle intrinsic to the grammar.
\end{remark}

% ------------------------------------------------------------
\section{Global Coverage and Density Exhaustion}
\label{sec:8a-global}

\begin{itemize}
    \item Iterated block covering argument
    \item Exhaustion of dyadic gaps
    \item No residual positive density set
    \item Candidate density-1 conclusion mechanism
\end{itemize}

% ============================================================

This establishes a strict subcritical regime for symbolic expansion.

% ============================================================
\clearpage
\bibliographystyle{unsrtnat}
\bibliography{src/latex/references}

\end{document}